% ==============================================================================
% 第 23 章：多输入同步处理 (Chapter 23: Simultaneous Processing)
% ==============================================================================

\chapter{多输入同步处理}
\label{sec:ch23_simultaneous}

某些函数必须接收两个（或更多）属于非平凡复杂数据定义的参数。如何设计此类函数，取决于这些参数之间的\textbf{结构对应关系}。大体可分为三类情形：
\begin{enumerate}
  \item \textbf{情形 1：被动参数。} 其中一个参数被当作原子值处理，仅对其余参数展开递归遍历。
  \item \textbf{情形 2：步调一致（齐头并进）。} 两个参数具有完全对称的相同结构，函数在两个参数上以同步锁步（lockstep）的方式并行前进。
  \item \textbf{情形 3：全部分支穷举。} 两个参数各自拥有多个分支，函数必须考虑两个参数数据分支的\textbf{笛卡尔积组合}。
\end{enumerate}

本章通过典型示例逐一阐述这三种情形，并给出扩展的设计配方。最后两节将探讨复合数据的相等性判定，以及一个完整的数据库（Database）关系代数查询引擎项目。

\section{同时处理两个列表：情形 1 (Processing Two Lists: Case 1)}
\label{sec:23.1_case1}

考虑如下函数签名、目的说明与函数头：
\begin{lstlisting}[language=BSL]
; [List-of Number] [List-of Number] -> [List-of Number]
; 将 front 末尾的 '() 替换为 end
(define (replace-eol-with front end)
  front)
\end{lstlisting}

签名表明该函数接收两个列表。我们来看设计配方在此如何展开。首先设计用例：
\begin{lstlisting}[language=BSL]
(check-expect (replace-eol-with '() '(a b)) '(a b))
(check-expect (replace-eol-with (cons 1 '()) '(a b)) (cons 1 '(a b)))
(check-expect (replace-eol-with (cons 2 (cons 1 '())) '(a b))
              (cons 2 (cons 1 '(a b))))
\end{lstlisting}

目的说明与示例清晰地表明：只要第二个参数 \code{end} 是一个合法的值，函数在整个执行过程中\textbf{完全不需要拆解它}。因此，模板只需针对第一个参数 \code{front} 建立标准的列表递归模板，而将 \code{end} 作为不变量随身传递：

\begin{lstlisting}[language=BSL]
(define (replace-eol-with front end)
  (cond
    [(empty? front) end]
    [else
     (cons (first front)
           (replace-eol-with (rest front) end))]))
\end{lstlisting}

这正是列表拼接操作（即内置函数 \code{append}）的标准实现！

\begin{exercise}{笛卡尔积 cross}{387}
设计函数 \code{cross}。该函数接收一个符号列表和一个数字列表，产生所有可能的有序符号-数字对（长度为 2 的列表）。例如当输入 \code{'(a b c)} 与 \code{'(1 2)} 时，预期结果为：
\begin{lstlisting}[language=BSL]
'((a 1) (a 2) (b 1) (b 2) (c 1) (c 2))
\end{lstlisting}
\end{exercise}

\section{同时处理两个列表：情形 2 (Processing Two Lists: Case 2)}
\label{sec:23.2_case2}

第 10.1 节介绍了根据工时计算周薪的 \code{wages*} 函数。当时假定所有工人的时薪统一。现实中，每位员工的时薪各不相同。我们设计更真实的进阶版本：函数同时接收\textbf{工时列表}与\textbf{对应时薪列表}。

\begin{lstlisting}[language=BSL]
; [List-of Number] [List-of Number] -> [List-of Number]
; 将 hours 与 wages/h 中的对应项逐一相乘
; 假设两个列表具有完全相同的长度
(define (wages*.v2 hours wages/h)
  '())
\end{lstlisting}

两个列表长度相等的先决条件意味着：\code{(empty? hours)} 为真当且仅当 \code{(empty? wages/h)} 为真；\code{(cons? hours)} 为真当且仅当 \code{(cons? wages/h)} 为真。因此，函数可以以\textbf{锁步（lockstep）}的方式同步拆解两个列表：

\begin{lstlisting}[language=BSL]
(define (wages*.v2 hours wages/h)
  (cond
    [(empty? hours) '()]
    [else
     (cons (weekly-wage (first hours) (first wages/h))
           (wages*.v2 (rest hours) (rest wages/h)))]))

(define (weekly-wage h r)
  (* h r))
\end{lstlisting}

\begin{exercise}{真实员工薪资计算}{388}
在现实中，\code{wages*.v2} 接收员工信息结构体列表与打卡记录结构体列表。员工结构体包含姓名、身份证号与时薪；打卡记录包含姓名与周工作小时数。请制定结构体定义与数据定义，并修改函数使其生成包含员工姓名与周薪的工资单结构体列表。
\end{exercise}

\begin{exercise}{列表打包函数 zip}{389}
设计函数 \code{zip}。它接收一个姓名列表（字符串）和一个电话号码列表（字符串），且两列表长度相等。函数将它们配对打包为电话簿记录列表：
\begin{lstlisting}[language=BSL]
(define-struct phone-record [name number])
; 一个 PhoneRecord 是：(make-phone-record String String)
\end{lstlisting}
\end{exercise}

\section{同时处理两个列表：情形 3 (Processing Two Lists: Case 3)}
\label{sec:23.3_case3}

当两个复杂参数之间既不存在从属关系，也不存在长度同步约束时，我们就进入了最一般的情形：

\begin{notebox}[样本问题]
给定符号列表 \code{los} 和自然数 \code{n}，设计函数 \code{list-pick} 提取 \code{los} 中索引为 \code{n} 的符号（索引从 0 开始计数）；若索引越界，则报错。
\end{notebox}

两个输入分别属于两个复杂数据定义：
\begin{itemize}
  \item \code{[List-of Symbol]} 有两个分支：\code{'()} 或 \code{(cons Symbol [List-of Symbol])}；
  \item \code{N}（自然数）有两个分支：\code{0} 或 \code{(add1 N)}。
\end{itemize}

两两组合构成了 $2 \times 2 = 4$ 种输入情形。我们用表格罗列这些组合：

\begin{center}
\begin{tabular}{|c|c|c|}
\hline
& \code{(= n 0)} & \code{(> n 0)} \\
\hline
\code{(empty? los)} & 列表为空，报错 \code{"list too short"} & 列表为空，报错 \code{"list too short"} \\
\hline
\code{(cons? los)} & 返回首元素 \code{(first los)} & 递归：\code{(list-pick (rest los) (sub1 n))} \\
\hline
\end{tabular}
\end{center}

据此直接导出的初始实现如图~\ref{fig:134} 所示。

\begin{figure}[htbp]
  \centering
\begin{lstlisting}[language=BSL, basicstyle=\scriptsize\ttfamily]
; [List-of Symbol] N -> Symbol
; extracts the nth symbol from l;
; signals an error if there is no such symbol
(define (list-pick l n)
  (cond
    [(and (= n 0) (empty? l))
     (error 'list-pick "list too short")]
    [(and (> n 0) (empty? l))
     (error 'list-pick "list too short")]
    [(and (= n 0) (cons? l)) (first l)]
    [(and (> n 0) (cons? l)) (list-pick (rest l) (sub1 n))]))
\end{lstlisting}
  \caption{基于全状态分支矩阵的列表索引检索 (Indexing into a list)}
  \label{fig:134}
\end{figure}

\begin{exercise}{树结构索引检索 tree-pick}{390}
设计函数 \code{tree-pick}。该函数接收一棵符号二叉树和一个由方向指令（\code{'left} 与 \code{'right}）组成的方向路径列表，沿着路径检索目标节点处的符号。
\end{exercise}

\section{简化函数结构}
\label{sec:23.4_simplification}

图~\ref{fig:134} 中的四个条件分支可以直接合并化简。前两个分支都在 \code{(empty? los)} 时抛出相同的错误，因此无论 \code{n} 取何值，只要列表为空就应立即报错。图~\ref{fig:135} 展示了精炼合并后的代码。

\begin{figure}[htbp]
  \centering
\begin{lstlisting}[language=BSL, basicstyle=\scriptsize\ttfamily]
; list-pick: [List-of Symbol] N[>= 0] -> Symbol
; determines the nth symbol from alos, counting from 0;
; signals an error if there is no nth symbol
(define (list-pick alos n)
  (cond
    [(empty? alos) (error 'list-pick "list too short")]
    [(= n 0) (first alos)]
    [(> n 0) (list-pick (rest alos) (sub1 n))]))
\end{lstlisting}
  \caption{化简后的列表索引函数 (Indexing into a list, simplified)}
  \label{fig:135}
\end{figure}

\begin{exercise}{用全分支分析法重构 replace-eol-with}{391}
使用 23.3 节的情形 3 分析策略重新设计 \code{replace-eol-with}。构建 $2 \times 2$ 的状态转移表格，然后化简得到最终函数，验证其与情形 1 版本的等价性。
\end{exercise}

\begin{exercise}{化简 tree-pick 函数}{392}
运用本节的方法，对练习 390 中的 \code{tree-pick} 函数实施化简。
\end{exercise}

\section{设计接收两个复杂输入的函数 }
\label{sec:23.5_recipe_two_inputs}

面对两个复杂参数时，我们可以将设计配方总结为如下矩阵设计法：

\begin{recipebox}[处理两个复杂输入的设计配方]
\textbf{1. 维度分析：} 分析输入参数的数据定义分支数。若参数 1 有 $M$ 个分支，参数 2 有 $N$ 个分支，建立一个 $M \times N$ 的二维表格。

\textbf{2. 逐格填充：} 为表格中的每一个单元格填入代表性的测试用例和期望结果。
\begin{itemize}
  \item 若某些行或列的结果完全独立于另一参数，则属于\textbf{情形 1}；
  \item 若仅对角线上的单元格合法且有意义，则属于\textbf{情形 2}；
  \item 若所有单元格均具有独特逻辑，则属于\textbf{情形 3}。
\end{itemize}

\textbf{3. 编写条件分支：} 将表格的每个有效单元格直接转译为 \code{cond} 中的一条条件分支。

\textbf{4. 化简合并：} 观察各分支的处理逻辑，利用布尔组合（\code{or}、\code{and}）合并抛出相同错误或执行相同操作的冗余分支。
\end{recipebox}

\section{练手习题：两个输入 }
\label{sec:23.6_finger_exercises}

\begin{exercise}{有限集合的并集}{393}
设计函数 \code{set-union}。该函数接收两个表示有限集合的无重复列表，计算它们的并集。
\end{exercise}

\begin{exercise}{合并两个有序列表 merge}{394}
设计函数 \code{merge}。该函数接收两个严格升序排列的数字列表，产生将二者合并后的单一升序排列列表。
\end{exercise}

\begin{exercise}{列表切片函数 take 与 drop}{395}
设计函数 \code{take}（获取列表的前 $n$ 项）与 \code{drop}（丢弃列表的前 $n$ 项）。
\end{exercise}

\begin{figure}[htbp]
  \centering
\begin{lstlisting}[language=BSL, basicstyle=\scriptsize\ttfamily]
; An HM-Word is a [List-of Letter or "_"]
; interpretation "_" represents a letter to be guessed

; HM-Word N -> String
; runs a simplistic hangman game, produces the current state
(define (play the-pick time-limit)
  (local ((define the-word  (explode the-pick))
          (define the-guess (make-list (length the-word) "_"))
          ; HM-Word -> HM-Word
          (define (do-nothing s) s)
          ; HM-Word KeyEvent -> HM-Word
          (define (checked-compare current-status ke)
            (if (member? ke LETTERS)
                (compare-word the-word current-status ke)
                current-status)))
    (implode
      (big-bang the-guess ; HM-Word
        [to-draw render-word]
        [on-tick do-nothing 1 time-limit]
        [on-key  checked-compare]))))

; HM-Word -> Image
(define (render-word w)
  (text (implode w) 22 "black"))
\end{lstlisting}
  \caption{简易猜词游戏 (A simple hangman game)}
  \label{fig:136}
\end{figure}

\begin{exercise}{猜词游戏 Hangman}{396}
图~\ref{fig:136} 展示了经典的猜词游戏 Hangman。请完善其中比对单词字母与猜测字符的核心判定逻辑。
\end{exercise}

\begin{exercise}{工时打卡记录匹配}{397}
工厂员工在上下班时打卡。设计函数计算所有员工的累计有效工时。
\end{exercise}

\begin{exercise}{线性组合计算}{398}
一个多元一次线性组合形如 $a_1 x_1 + a_2 x_2 + \dots + a_n x_n$。设计函数接收系数列表与变量值列表，计算其代数和。
\end{exercise}

\begin{exercise}{抽奖礼物互赠配对}{399}
五位朋友决定举行秘密礼物互赠抽奖。设计函数产生一个不包含任何人抽到自己（即无不动点的错排，derangement）的赠礼配对列表。
\end{exercise}

\begin{exercise}{DNA 前缀匹配 DNAprefix}{400}
设计函数 \code{DNAprefix}，判断一个核苷酸碱基序列是否是另一个较长 DNA 序列的严格前缀。
\end{exercise}

\begin{exercise}{S-表达式相等性判定 sexp=?}{401}
设计通用谓词函数 \code{sexp=?}，判断两个任意构造的 S-表达式是否完全结构等价。
\end{exercise}

\begin{exercise}{复合数据的相等性探讨}{402}
探讨判断两个复杂数据结构相等的本质难点：为什么我们无法为通用函数本身编写一个 \code{fun=?} 判定函数？（停机问题与哥德尔不完备定理的直观隐喻）。
\end{exercise}

\section{项目：数据库 }
\label{sec:23.7_database}

关系数据库管理系统（RDBMS）是现代软件工业的基石。在本质上，关系数据库将数据存储为二维表格（Table），并通过关系代数（如投影、筛选、笛卡尔积、连接）进行结构化查询。

图~\ref{fig:137} 展示了两个表格实例：学生表与出勤表。

\begin{figure}[htbp]
  \centering
\begin{stepperbox}[unbreakable]
  \centering\scriptsize
  \begin{tabular}[t]{@{}lll@{}}
    \toprule
    \textbf{Name} & \textbf{Age} & \textbf{Present} \\
    String & Integer & Boolean \\
    \midrule
    \texttt{"Alice"} & \texttt{35} & \texttt{\#true} \\
    \texttt{"Bob"}   & \texttt{25} & \texttt{\#false} \\
    \texttt{"Carol"} & \texttt{30} & \texttt{\#true} \\
    \texttt{"Dave"}  & \texttt{32} & \texttt{\#false} \\
    \bottomrule
  \end{tabular}%
  \hspace{1.2em}%
  \begin{tabular}[t]{@{}ll@{}}
    \toprule
    \textbf{Present} & \textbf{Description} \\
    Boolean & String \\
    \midrule
    \texttt{\#true}  & \texttt{"presence"} \\
    \texttt{\#false} & \texttt{"absence"} \\
    \bottomrule
  \end{tabular}
\end{stepperbox}
  \caption{作为二维数据表的数据库 (Databases as tables)}
  \label{fig:137}
\end{figure}

图~\ref{fig:138} 展示了在 ISL+ 中对数据库及其模式（Schema）的严谨数据定义。

\begin{figure}[htbp]
  \centering
\begin{stepperbox}[unbreakable]
\noindent\begin{minipage}[t]{0.48\linewidth}\vspace{0pt}
\begin{Verbatim}[fontsize=\fontsize{6.0pt}{7.2pt}\selectfont]
(define school-schema
  `(("Name"    ,string?)
    ("Age"     ,integer?)
    ("Present" ,boolean?)))

(define school-content
  `(("Alice" 35 #true)
    ("Bob"   25 #false)
    ("Carol" 30 #true)
    ("Dave"  32 #false)))

(define school-db
  (make-db school-schema
           school-content))
\end{Verbatim}
\end{minipage}\hfill
\noindent\begin{minipage}[t]{0.48\linewidth}\vspace{0pt}
\begin{Verbatim}[fontsize=\fontsize{6.0pt}{7.2pt}\selectfont]
(define presence-schema
  `(("Present"     ,boolean?)
    ("Description" ,string?)))

(define presence-content
  `((#true  "presence")
    (#false "absence")))

(define presence-db
  (make-db presence-schema
           presence-content))
\end{Verbatim}
\end{minipage}
\end{stepperbox}
  \caption{数据库在 ISL+ 中的数据表示 (Databases as ISL+ data)}
  \label{fig:138}
\end{figure}

\begin{exercise}{规范定义 Spec}{403}
模式规范定义将列名（Label）与该列值的校验谓词（Predicate）关联起来。为 \code{Spec} 建立结构体定义与验证函数。
\end{exercise}

\begin{exercise}{双列表断言函数 andmap2}{404}
设计函数 \code{andmap2}。它接收一个二元谓词函数与两个等长列表，仅当对应位置的元素对全部满足该谓词时才返回 \code{\#true}。
\end{exercise}

图~\ref{fig:139} 与图~\ref{fig:140} 给出了数据库列投影运算的设计模板与表达式外提（hoisting）优化技术。

\begin{figure}[htbp]
  \centering
\begin{lstlisting}[language=BSL, basicstyle=\scriptsize\ttfamily]
(define (integrity-check.v3 db)
  (local ((define schema  (db-schema db))
          (define content (db-content db))
          (define width   (length schema))
          ; Row -> Boolean
          ; does row satisfy (I1) and (I2)
          (define (row-integrity-check row)
            (and (= (length row) width)
                 (andmap (lambda (s c) [(second s) c])
                         schema
                         row))))
    (andmap row-integrity-check content)))
\end{lstlisting}
  \caption{系统的表达式外提优化结果 (The result of systematic expression hoisting)}
  \label{fig:139}
\end{figure}

\begin{figure}[htbp]
  \centering
\begin{lstlisting}[language=BSL, basicstyle=\scriptsize\ttfamily]
(define (project db labels)
  (local ((define schema  (db-schema db))
          (define content (db-content db))
          ; Spec -> Boolean
          ; does this spec belong to the new schema
          (define (keep? c) ...)
          ; Row -> Row
          ; retains those columns whose name is in labels
          (define (row-project row) ...))
    (make-db (filter keep? schema)
             (map row-project content))))
\end{lstlisting}
  \caption{数据库投影项目模板 (A template for project)}
  \label{fig:140}
\end{figure}

\begin{exercise}{行过滤函数 row-filter}{405}
设计函数 \code{row-filter}。该函数根据给定的谓词条件筛选数据表中符合条件的行记录。
\end{exercise}

图~\ref{fig:141} 与图~\ref{fig:142} 说明了投影运算（\code{project}）的运作全貌。

\begin{figure}[htbp]
  \centering
\begin{lstlisting}[language=BSL, basicstyle=\scriptsize\ttfamily]
(define (project.v1 db labels)
  (local ((define schema  (db-schema db))
          (define content (db-content db))

          ; Spec -> Boolean
          ; does this column belong to the new schema
          (define (keep? c)
            (member? (first c) labels))

          ; Row -> Row
          ; retains those columns whose name is in labels
          (define (row-project row)
            (row-filter row (map first schema)))

          ; Row [List-of Label] -> Row
          ; retains those cells whose name is in labels
          (define (row-filter row names)
            (cond
              [(empty? names) '()]
              [else
               (if (member? (first names) labels)
                   (cons (first row)
                         (row-filter (rest row) (rest names)))
                   (row-filter (rest row) (rest names)))])))
    (make-db (filter keep? schema)
             (map row-project content))))
\end{lstlisting}
  \caption{数据库列投影示意图 (Database projection)}
  \label{fig:141}
\end{figure}

\begin{figure}[htbp]
  \centering
\begin{lstlisting}[language=BSL, basicstyle=\scriptsize\ttfamily]
(define (project db labels)
  (local ((define schema  (db-schema db))
          (define content (db-content db))

          ; Spec -> Boolean
          ; does this column belong to the new schema
          (define (keep? c)
            (member? (first c) labels))

          ; Row -> Row
          ; retains those columns whose name is in labels
          (define (row-project row)
            (foldr (lambda (cell m c) (if m (cons cell c) c))
                   '()
                   row
                   mask))
          (define mask (map keep? schema)))
    (make-db (filter keep? schema)
             (map row-project content))))
\end{lstlisting}
  \caption{基于模式投影的完整流程 (Database projection)}
  \label{fig:142}
\end{figure}

\begin{exercise}{投影标签优化}{406}
优化投影算法，避免在处理每一行时重复计算保留列的索引位置。
\end{exercise}

\begin{exercise}{使用 foldr 重构行过滤}{407}
使用 \code{foldr} 重新设计 \code{row-filter}。
\end{exercise}

\begin{exercise}{数据库投影查询 select}{408}
设计函数 \code{select}。它接收一个数据库和一组目标列名列表，产生仅包含这些列的新数据库。
\end{exercise}

\begin{exercise}{数据库行排序 reorder}{409}
设计函数 \code{reorder}，按照指定列的值对数据库中的行记录进行排序。
\end{exercise}

\begin{exercise}{数据库并集 db-union}{410}
设计函数 \code{db-union}，计算具有相同模式的两个数据库表的行并集（自动去重）。
\end{exercise}

\begin{exercise}{数据库自然连接 join}{411}
设计数据库关系代数中最核心的运算：自然连接函数 \code{join}。它接收具有公共列的两个数据表，根据公共列上取值相等的条件将它们拼接为包含更丰富信息的新表。
\end{exercise}
